% !TEX TS-program = pdflatex
% !TeX program = pdflatex
% !TEX encoding = UTF-8
% !TEX spellcheck = en_US

\documentclass[xcolor=table]{beamer}


%\usepackage{fullpage}
%\usepackage[left=2.8cm,right=2.2cm,top=2 cm,bottom=2 cm]{geometry}
\setbeamersize{text margin left=10pt,text margin right=10pt}



%\usepackage[T1]{fontenc}
\usepackage{textcomp}
\usepackage[utf8]{inputenc}
%\usepackage[french,english]{babel}
\usepackage{arabtex}
\usepackage{txfonts}
\usepackage{tipa}
\usepackage[]{graphicx}
\usepackage{multirow}
\usepackage{hyperref}
\usepackage{colortbl}
\usepackage{tabularray}
\usepackage{listingsutf8}
\usepackage{wrapfig}
\usepackage{multicol}
\usepackage[export]{adjustbox} %for images in table, also for frame
\usepackage[many]{tcolorbox}
\usepackage{wasysym}
\usepackage[lined]{algorithm2e}
\usepackage{alltt} %verbatim with commands
\usepackage{longtable}
\usepackage{tabu}

\usepackage{CJKutf8}

\usepackage{amsmath,amssymb, amsfonts} 

%\usepackage{newtxtext,newtxmath}


%\usepackage{acolor}


\definecolor{lightblue}{HTML}{D0D2FF}
\definecolor{lightyellow}{HTML}{FFFFAA}
\definecolor{darkblue}{HTML}{0000BB}
\definecolor{olivegreen}{HTML}{006600}
\definecolor{darkgreen}{HTML}{008B45} %009B55
\definecolor{violet}{HTML}{6600CC}
\definecolor{deeppink}{HTML}{FF1493}
\definecolor{orangey}{HTML}{FFBB00}


\makeatletter
\def\beamer@calltheme#1#2#3{%
	\def\beamer@themelist{#2}
	\@for\beamer@themename:=\beamer@themelist\do
	{\IfFileExists{\beamer@themelocation/#3\beamer@themename.sty}{%
	\usepackage[{#1}]{\beamer@themelocation/#3\beamer@themename}}{\usepackage[{#1}]{#3\beamer@themename}}
		
	}%
}

\def\usefolder#1{
	\def\beamer@themelocation{#1}
}
\def\beamer@themelocation{}
\makeatother

\usefolder{../../extra/beamer-themes}

\usetheme{Karim} % Antibes Boadilla Warsaw Karim

\beamertemplatenavigationsymbolsempty

\let\oldcite\cite
\renewcommand{\cite}[1]{{\bfseries\color{blue}\oldcite{#1}}}

\tcbuselibrary{listings}


%\renewcommand{\baselinestretch}{1.5}

\def\supit#1{\raisebox{0.8ex}{\small\it #1}\hspace{0.05em}}

\AtBeginSection{%
	\begin{frame}
		\sectionpage
	\end{frame}
}

\newcommand{\rottext}[2]{%
	\rotatebox{90}{%
	\begin{minipage}{#1}%
		\raggedleft#2%
	\end{minipage}%
	}%
}


\institute{ %
Laboratoire de la Communication dans les Systèmes Informatiques (LCSI)
\par
École  nationale Supérieure d'Informatique (ESI), Algiers, Algeria
}
\author[Abdelkrime Aries (2023/2024)] %
{Abdelkrime Aries}
%\titlegraphic{\includegraphics[height=1cm]{../img/esi-logo.png}%\hspace*{4.75cm}~


\date{Academic year: 2023/2024} %\today

\titlegraphic{%
	\includegraphics[height=1cm]{../../img/misc/esi-logo.png}%
	\hspace{2cm}%
	\includegraphics[height=1cm]{../../img/misc/lcsi-logo.png}%
	\hspace{2cm}%
	\includegraphics[height=1.5cm]{../../img/misc/esi.nlp-logo.png}%
}

%\setbeamertemplate{headline}{}

\newcommand{\kurl}[1]{{\scriptsize\bfseries\color{orangey}\url{#1}}}

\newcommand{\keyword}[1]{\textcolor{red}{\bfseries\itshape #1}}
\newcommand{\expword}[1]{\textcolor{olivegreen}{#1}}
\newcommand{\optword}[1]{\textcolor{violet}{\bfseries #1}}

\makeatletter
\newcommand\mysphere{%
	\parbox[t]{10pt}{\raisebox{0.2pt}{\beamer@usesphere{item projected}{bigsphere}}}}
\makeatother

%\let\oldtabular\tabular
%\let\endoldtabular\endtabular
%\renewenvironment{tabular}{\rowcolors{2}{white}{lightblue}\oldtabular\rowcolor{blue}}{\endoldtabular}


%\NoAutoSpacing %french autospacing after ":"

\def\graphpath{}

\newcommand{\changegraphpath}[1]{\def\graphpath{#1}}


\newcommand{\vgraphpage}[2][.84\textheight]{%
%	\begin{center}%
		\includegraphics[height=#1]{\graphpath #2}%
%	\end{center}%
}

\newcommand{\hgraphpage}[2][\textwidth]{%
%	\begin{center}%
		\includegraphics[width=#1]{\graphpath #2}%
%	\end{center}%
}

\newcommand{\graphpage}[2][]{%
	\includegraphics[#1]{\graphpath #2}%
}

\bibliographystyle{apalike}

\newcommand{\insertbibliography}[2]{
	\appendix
	\section*{Bibliography}
	\nocite{#2}
%	\makeatletter % to change template
%	\setbeamertemplate{headline}[default] % not mandatory, but I though it was better to set it blank
%	\def\beamer@entrycode{\vspace*{-\headheight}} % here is the part we are interested in :)
%	\makeatother
	\begin{multicols*}{2}[\frametitle{\insertsection} \usebeamertemplate*{frametitle}]%\usebeamertemplate*{frametitle}\frametitle{Références}
		\tiny
		\bibliography{#1}
	\end{multicols*}
}

\definecolor{my-grey}{RGB}{233, 233, 233}

\newcommand{\insertlicence}{
	\begin{frame}[plain]
	\frametitle{License: CC-BY 4.0}

	\begin{tcolorbox}[colback=cyan,
		colframe=cyan,  
		arc=0pt,outer arc=0pt,
		valign=top, 
		halign=center,
		width=\textwidth]
		
		\includegraphics[width=.5cm]{../../img/licence/cc_icon_white_x2.png}
		\includegraphics[width=.5cm]{../../img/licence/attribution_icon_white_x2.png}
		
		\color{white}
		\bfseries Attribution 4.0 International (CC BY 4.0) \\
		\tiny \url{https://creativecommons.org/licenses/by/4.0/deed.en}
		
	\end{tcolorbox}\vspace{-.5cm}
	\begin{tcolorbox}[colback=my-grey,
		colframe=my-grey,  
		center, arc=0pt,outer arc=0pt,
		valign=top, 
		halign=left,
		width=\textwidth]
		
		\tiny
		
		\begin{center}
			\bfseries\Large
			You are free to:
		\end{center}
		
		\begin{minipage}{0.83\textwidth}
			\begin{itemize}
				\item[] \textbf{Share} — copy and redistribute the material in any medium or format
				\item[] \textbf{Adapt} — remix, transform, and build upon the material
				for any purpose, even commercially.
			\end{itemize}
		
			\vspace{12pt}
		
			\textit{The licensor cannot revoke these freedoms as long as you follow the license terms.}
		\end{minipage}
		\begin{minipage}{0.15\textwidth}
			\includegraphics[width=\textwidth]{../../img/licence/FreeCulturalWorks_seal_x2.jpg}
		\end{minipage}
	
		
		\begin{center}
			\bfseries\Large
			Under the following terms:
		\end{center}
		
		\begin{itemize}
			\item[] \textbf{Attribution} — You must give appropriate credit, provide a link to the license, and indicate if changes were made. You may do so in any reasonable manner, but not in any way that suggests the licensor endorses you or your use.
			\item[] \textbf{No additional restrictions} — You may not apply legal terms or technological measures that legally restrict others from doing anything the license permits.
		\end{itemize}
		
	\end{tcolorbox}
	
%	\begin{center}
%		\bfseries Attribution 4.0 International (CC BY 4.0)
%		\url{https://creativecommons.org/licenses/by/4.0/deed.fr}
%	\end{center}

%	\tiny
%
%	Vous êtes autorisé à : 
%	\begin{itemize}
%		\item \textbf{Partager} — copier, distribuer et communiquer le matériel par tous moyens et sous tous formats
%		\item \textbf{Adapter} — remixer, transformer et créer à partir du matériel
%	\end{itemize}
%	
%	Selon les conditions suivantes : 
%	\begin{itemize}
%		\item \textbf{Attribution} — Vous devez créditer l'Œuvre, intégrer un lien vers la licence et indiquer si des modifications ont été effectuées à l'Oeuvre. Vous devez indiquer ces informations par tous les moyens raisonnables, sans toutefois suggérer que l'Offrant vous soutient ou soutient la façon dont vous avez utilisé son Oeuvre.
%		\item \textbf{Pas d'Utilisation Commerciale} — Vous n'êtes pas autorisé à faire un usage commercial de cette Oeuvre, tout ou partie du matériel la composant. 
%		\item \textbf{Pas de restrictions complémentaires} — Vous n'êtes pas autorisé à appliquer des conditions légales ou des mesures techniques qui restreindraient légalement autrui à utiliser l'Oeuvre dans les conditions décrites par la licence.
%	\end{itemize}

	\end{frame}
}

\settowidth{\leftmargini}{\usebeamertemplate{itemize item}}
\addtolength{\leftmargini}{\labelsep}

\AtBeginDocument{
	\newcolumntype{L}[2]{>{\vbox to #2\bgroup\vfill\flushleft}p{#1}<{\egroup}} 
	
	\begin{frame}[plain]
		\maketitle
	\end{frame}

	\insertlicence
}


% needs etoolbox; to break links after -
\appto\UrlBreaks{\do\-}


\makeatletter
\newcommand{\xRightarrow}[2][]{\ext@arrow 0359\Rightarrowfill@{#1}{#2}}
\makeatother


\usefonttheme{structurebold}
%\usefonttheme{professionalfonts}


% figure caption
\setlength\abovecaptionskip{2pt}

% equations
\AtBeginDocument{
\abovedisplayskip=0pt
\abovedisplayshortskip=0pt
\belowdisplayskip=0pt
\belowdisplayshortskip=0pt
}


\title[ESI - NLP: 07- Word semantics]%
{Natural Language Processing\\Chapter 07\\Word semantics} 

\changegraphpath{../img/word-sem/}

\begin{document}
	
\begin{frame}
\frametitle{Natural Language Processing}
\framesubtitle{Word semantics: Introduction}

\begin{exampleblock}{Example of polysemy}
	\begin{center}
		\Large\bfseries
%		\begin{itemize}
			The hotel has a big \expword{wing}.
			 
			That bird broke its \expword{wing}.
			
			They attacked the enemy's right \expword{wing}.
			
			I ate just a \expword{wing}.
%		\end{itemize}
	\end{center}
\end{exampleblock}

\begin{itemize}
	\item Does the word ``\expword{wing}" mean the same thing in all four sentences?
	\item How can we know the meaning of a word in this case?
	\item The meanings of the word ``\expword{wing}": \url{http://wordnetweb.princeton.edu/perl/webwn?s=wing}
\end{itemize}

\end{frame}

\begin{frame}
\frametitle{Natural Language Processing}
\framesubtitle{Word semantics: Plan}

\begin{multicols}{2}
%	\small
\tableofcontents
\end{multicols}
\end{frame}

%===================================================================================
\section{Lexical databases}
%===================================================================================

\begin{frame}
	\frametitle{Word semantics}
	\framesubtitle{Lexical databases}

	\begin{itemize}
		\item A database containing a language's vocabulary
		\item Word information: PoS, lemma, frequency, etc.
	\end{itemize}

	\begin{figure}
		\centering 
		\hgraphpage[.5\textwidth]{exp-bd-lex.pdf}
		\caption{Example of categories and their relations in a lexical base \cite{2019-white-al}}
	\end{figure}
	
\end{frame}

\subsection{Semantic relations}

\begin{frame}
	\frametitle{Word semantics: Lexical DBs}
	\framesubtitle{Semantic relations (Reminder)}

	\begin{itemize}
		\item \optword{Synonymy}: having similar meanings given a context.
		\item \optword{Antonymy}: having opposite meanings given a context.
		\item Taxonomic (classification) relations
		\begin{itemize}
			\item \optword{Hyponymy}: be more specific than another meaning. It results in a \keyword{IS-A} relation. E.g. \expword{car IS-A vehicle}.
			\item \optword{Hyperonymy}: to be more generic than another meaning.
			\item \optword{Meronymy}: to be part of a thing. E.g. \expword{wheel is a meronym of car; car is the holonym of wheel}
		\end{itemize}
	\end{itemize}
	
\end{frame}

\subsection{Wordnet}

\begin{frame}
\frametitle{Word semantics: Lexical DBs}
\framesubtitle{Wordnet}

\begin{exampleblock}{Example of WordNet : \small\url{http://wordnetweb.princeton.edu/perl/webwn}}
	\fontsize{6}{6}\selectfont
	%		\begin{alltt}
	{\small\bfseries Noun}
	\begin{itemize}\setlength\itemsep{0pt}
		\item \textcolor{blue}{\underline{S:}} \textcolor{red}{(n)} \textbf{reason}, \textcolor{blue}{\underline{ground}} (a rational motive for a belief or action) \textit{"the reason that war was declared"; "the grounds for their declaration"}
		\item \textcolor{blue}{\underline{S:}} \textcolor{red}{(n)} \textbf{reason} (an explanation of the cause of some phenomenon) \textit{"the reason a steady state was never reached was that the back pressure built up too slowly"}
		\item \textcolor{blue}{\underline{S:}} \textcolor{red}{(n)} \textbf{reason}, \textcolor{blue}{\underline{understanding}}, \textcolor{blue}{\underline{intellect}} (the capacity for rational thought or inference or discrimination) \textit{"we are told that man is endowed with reason and capable of distinguishing good from evil"}
		\item \textcolor{blue}{\underline{S:}} \textcolor{red}{(n)} \textcolor{blue}{\underline{rationality}}, \textbf{reason}, \textcolor{blue}{\underline{reasonableness}} (the state of having good sense and sound judgment) \textit{"his rationality may have been impaired"; "he had to rely less on reason than on rousing their emotions"}
		\item \textcolor{blue}{\underline{S:}} \textcolor{red}{(n)} \textcolor{blue}{\underline{cause}}, \textbf{reason}, \textcolor{blue}{\underline{grounds}} (a justification for something existing or happening) \textit{"he had no cause to complain"; "they had good reason to rejoice"}
		\item \textcolor{blue}{\underline{S:}} \textcolor{red}{(n)} \textbf{reason} (a fact that logically justifies some premise or conclusion) \textit{"there is reason to believe he is lying"}
	\end{itemize}
	
	{\small\bfseries Verb}
	\begin{itemize}\setlength\itemsep{0pt}
		\item \textcolor{blue}{\underline{S:}} \textcolor{red}{(v)} \textbf{reason}, \textcolor{blue}{\underline{reason out}}, \textcolor{blue}{\underline{conclude}} (decide by reasoning; draw or come to a conclusion) \textit{"We reasoned that it was cheaper to rent than to buy a house"}
		\item \textcolor{blue}{\underline{S:}} \textcolor{red}{(v)} \textcolor{blue}{\underline{argue}}, \textbf{reason} (present reasons and arguments)
		\item \textcolor{blue}{\underline{S:}} \textcolor{red}{(v)} \textbf{reason} (think logically) \textit{"The children must learn to reason"}
	\end{itemize}
\end{exampleblock}
	
\end{frame}

\begin{frame}
\frametitle{Word semantics: Lexical DBs}
\framesubtitle{Wordnet: Description}
	
\begin{itemize}
	\item English lexical database \cite{1995-miller}.
	\item Three parts: (1) nouns (2) verbs (3) adjectives and adverbs.
	\item Each sens is represented by an identifier (\keyword{synset}).
	\begin{itemize}
		\item \keyword{synset} (Synonyms set): combines words with the same meaning. 
		\item E.g. \expword{05659525: reason\#3, understanding\#4, intellect\#2}
	\end{itemize}
	\item Each sens is defined by a glossary (\keyword{gloss}).
	\begin{itemize}
		\item E.g. \expword{05659525: (the capacity for rational thought or inference or discrimination) "we are told that man is endowed with reason and capable of distinguishing good from evil"}
	\end{itemize}
	\item Each sens has a lexicographic category (\keyword{supersense}).
	\begin{itemize}
		\item E.g. \expword{05659525: noun.cognition}
	\end{itemize}
	\item WordNet represents semantic relations between senses.
\end{itemize}
	
\end{frame}

\begin{frame}
\frametitle{Word semantics: Lexical DBs}
\framesubtitle{Wordnet: Lexicographic categories (supersense)}
	
\begin{block}{WordNet nouns lexicographic categories \cite{2019-jurafsky-martin}}
	\fontsize{8}{12}\selectfont\bfseries
	\centering
	\begin{tblr}{
		colspec = {llllll},
		row{odd} = {lightblue},
		row{even} = {lightyellow},
		row{1} = {darkblue},
		rowsep=1pt,
		colsep=3pt,
	} 
		\textcolor{white}{Category} & \textcolor{white}{Example} & \textcolor{white}{Category} & \textcolor{white}{Example} &\textcolor{white}{Category} & \textcolor{white}{Example} \\
%		\hline
		ACT & service & GROUP & place & PLANT & tree \\
		ANIMAL &  dog & LOCATION & area & POSSESSION & price \\
		ARTIFACT & car & MOTIVE & reason & PROCESS & process \\
		ATTRIBUTE & quality & NATURAL EVENT & experience & QUANTITY & amount \\
		BODY & hair & NATURAL OBJECT & flower & RELATION & portion \\
		COGNITION & way & OTHER & stuff & SHAPE & square\\
		COMMUNICATION & review & PERSON & people & STATE & pain\\
		FEELING & discomfort & PHENOMENON & result & SUBSTANCE & oil \\
		FOOD & food & & & TIME & day\\
%		\hline\hline
	\end{tblr}
\end{block}
	
\end{frame}

\begin{frame}
\frametitle{Word semantics: Lexical DBs}
\framesubtitle{Wordnet: Semantic relations (Nouns)}

\begin{block}{Some nouns relations \cite{2019-jurafsky-martin}}
	\fontsize{7}{14}\selectfont\bfseries\centering
	\begin{tblr}{
			colspec = {lll},
			row{odd} = {lightblue},
			row{even} = {lightyellow},
			row{1} = {darkblue},
			rowsep=1pt,
			colsep=4pt,
		} 
		\textcolor{white}{Relation} & \textcolor{white}{Definition} & \textcolor{white}{Example} \\
		Hypernym & from a specific concept to a more generic one & \expword{breakfast\textsuperscript{1} \textrightarrow\ meal\textsuperscript{1} }\\
		Hyponym & from a generic concept to a more specific one  & \expword{meal\textsuperscript{1} \textrightarrow\ lunch\textsuperscript{1}} \\
		Instance Hypernym & from an instance to its concept & \expword{Austen\textsuperscript{1} \textrightarrow\ author\textsuperscript{1}} \\
		Instance Hyponym & from an instance to its concept & \expword{composer\textsuperscript{1} \textrightarrow\ Bach\textsuperscript{1}} \\
		Part Meronym & from an entire concept to its part & \expword{table\textsuperscript{2} \textrightarrow\ leg\textsuperscript{3}} \\
		Part Holonym & from a part to its entirety & \expword{course\textsuperscript{7} \textrightarrow\ meal\textsuperscript{1}} \\
		Antonym & from a concept to its semantic opposition & \expword{leader\textsuperscript{1} $ \leftrightarrow $ follower\textsuperscript{1}}\\
		Derivation & from a word to another having the same root & \expword{destruction\textsuperscript{1} $ \leftrightarrow $ destroy\textsuperscript{1}} \\
	\end{tblr}
\end{block}
	
\end{frame}

\begin{frame}
	\frametitle{Word semantics: Lexical DBs}
	\framesubtitle{Wordnet: Semantic relations (Verbs)}
	
	\begin{block}{Some verbs relations \cite{2019-jurafsky-martin}}
		\fontsize{7}{14}\selectfont\bfseries\centering
		\begin{tblr}{
				colspec = {lll},
				row{odd} = {lightblue},
				row{even} = {lightyellow},
				row{1} = {darkblue},
				rowsep=1pt,
				colsep=4pt,
			} 
			\textcolor{white}{Relation} & \textcolor{white}{Definition} & \textcolor{white}{Example} \\
			Hypernym & from a specific event to a more generic one & \expword{fly\textsuperscript{9} \textrightarrow\ travel\textsuperscript{5}} \\
			Troponym & from a generic event to a more specific one & \expword{walk\textsuperscript{1} \textrightarrow\ stroll\textsuperscript{1}} \\
			Entails & from an event to another that it entails & \expword{snore\textsuperscript{1} \textrightarrow\ sleep\textsuperscript{1}} \\ 
			Antonym & from an event to its semantic opposition & \expword{increase\textsuperscript{1} $ \leftrightarrow $ decrease\textsuperscript{1}} \\
%			\hline\hline
		\end{tblr}
	\end{block}
	
\end{frame}

\begin{frame}
\frametitle{Word semantics: Lexical DBs}
\framesubtitle{Wordnet: Some resources}
	
\begin{itemize}
	\item Some APIs
	\begin{itemize}
		\item NLTK (Python): \url{https://www.nltk.org/howto/wordnet.html}
		\item JWI (Java): \url{http://projects.csail.mit.edu/jwi}
		\item Wordnet (Ruby): \url{https://github.com/wordnet/wordnet}
		\item OpenNlp (C\#): \url{https://github.com/AlexPoint/OpenNlp}
	\end{itemize}
	\item Other languages
	\begin{itemize}
		\item Global WordNet Association: \url{http://globalwordnet.org/resources/wordnets-in-the-world/}
		\item Open Multilingual Wordnet: \url{http://compling.hss.ntu.edu.sg/omw/}
	\end{itemize}
\end{itemize}
	
\end{frame}

\begin{frame}
	\frametitle{Word semantics: Lexical DBs}
	\framesubtitle{Wordnet: Similarity measures}
	
	\begin{itemize}
		\item Let us define some measures over concepts $c_i$ and $c_j$:
		\begin{itemize}
			\item $len(c_i, c_j)$: the length of the shortest IS-A path from $c_i$ to $c_j$
			\item $lso(c_i, c_j)$: the lowest common subsumer (most specific ancestor node)
			\item $depth(c_i)$: the length of the path from $c_i$ to the root entity
			\item $deep\_max$: the max $depth(c_i)$ of the taxonomy
			\item $IC(c_i)$: the probability of the concept given a corpus
		\end{itemize}
		\item Similarity measures
		\begin{itemize}
			\item \optword{Path similarity}: $sim_{path}(c_i, c_j) = \frac{1}{1+len(c_i, c_j)}$
			\item \optword{Leacock Chodorow Similarity}: $sim_{lch}(c_i, c_j) = -\log\frac{len(c_i, c_j)}{2*deep\_max}$
			\item \optword{Wu-Palmer Similarity}: $sim_{wup}(c_i, c_j) = \frac{2*depth(lso(c_i, c_j))}{len(c_i, c_j)+2*depth(lso(c_i, c_j))}$
			\item \optword{Lin Similarity}: $sim_{lin}(c_i, c_j) = \frac{2*IC(lso(c_i, c_j))}{IC(c_i)+IC(c_j)}$ 
			\item ...
		\end{itemize}
	\end{itemize}
	
\end{frame}

\begin{frame}
	\frametitle{Word semantics: Lexical DBs}
	\framesubtitle{Wordnet: Exercise}
	
	\hgraphpage{wn-exp.pdf}
	
	\begin{itemize}
		\item Calculate path similarity between all concepts if applicable
	\end{itemize}
	
\end{frame}

\subsection{Other resources}

\begin{frame}
\frametitle{Word semantics: Lexical DBs}
\framesubtitle{Other resources: Other Lexical DBs}
	
\begin{itemize}
	\item \optword{FrameNet} 
	\begin{itemize}
		\item Based on frame semantics theory.
		\item A frame can be an event, a relation or an entity with its participants.
		\item E.g. \expword{The concept ``Cook" imply a cook (person), prepared food, a container and a heating instrument}.
		\item Each frame is activated by a set of \keyword{lexical units}. E.g. \expword{bake, cook, fry, grill, prepare, roast}.
	\end{itemize}
	\item \optword{VerbNet}:
	\begin{itemize}
		\item A lexical base for verbs.
		\item It includes 30 principal thematic roles.
		\item The verbs are organized in classes.
	\end{itemize}
	\item \textbf{We will see these two DBs again in the next chapter.}
\end{itemize}
	
\end{frame}

\begin{frame}
\frametitle{Word semantics: Lexical DBs}
\framesubtitle{Other resources: BabelNet}

\begin{itemize}
	\item Multilingual sematic network
	\item \url{https://babelnet.org/}
\end{itemize}

\begin{figure}
	\hgraphpage{babelnet.pdf}
	\caption{BabelNet's structure \cite{2012-navigli-ponzetto}}
\end{figure}
	
\end{frame}

%===================================================================================
\section{Vector representation of words}
%===================================================================================

\begin{frame}
\frametitle{Word semantics}
\framesubtitle{Vector representation of words}

\begin{itemize}
	\item A word can be represented using \keyword{One-Hot} encoding:
	\begin{itemize}
		\item It does not represent semantic relations between words: similarity (E.g. \expword{cat, dog}) and proximity (E.g. \expword{coffee, cup}).
		\item How to represent a document/sentence using this representation?
	\end{itemize}

	\item \optword{Term-document}:
	\begin{itemize}
		\item A term is represented using the documents containing it (or the inverse).
		\item E.g. \expword{Term frequency - inverse document frequency (Tf-Idf)}.
	\end{itemize}

	\item \optword{Term-term}:
	\begin{itemize}
		\item A term is represented using other terms.
	\end{itemize}

	\item \optword{Term-concept-document}:
	\begin{itemize}
		\item Terms and documents are represented using a vector of concepts.
		\item E.g. \expword{Latent Semantic Analysis (LSA)}.
		\item Welcome to \keyword{Word embedding}!
	\end{itemize}

\end{itemize}

\end{frame}


\subsection{TF-IDF}

\begin{frame}
\frametitle{Word semantics: Vector representation}
\framesubtitle{TF-IDF}

\begin{itemize}
	\item A document/sentence sens can be represented using the words it contains. 
	\item Word's frequency in a document/sentence is called \keyword{term frequency (TF)}.
	\item There are words that are repeated a lot, which have no added meaning (E.g. \expword{
	Prepositions}).
	\item Words that appear in multiple documents are less important.
\end{itemize}

\begin{block}{TF-IDF calculation}
	\[
	TF_d(t) =  |\{t_i \in d / t_i = t\}|
	\hskip2cm 
	DF_D(t) = |\{d \in D / t \in d\}|
	\]
	\[IDF_D(t) = \log_{10} \left( \frac{|\{d \in D\}|}{DF_D(t)} \right)\]
	\[TF\text{-}IDF_{d, D}(t) = TF_d(t) * IDF_D(t)\]
\end{block}

\end{frame}


\begin{frame}
\frametitle{Word semantics: Vector representation}
\framesubtitle{TF-IDF: Example of TF representation}

\begin{exampleblock}{Example of some sentences}
	\begin{itemize}
		\item S1: a computer can help you
		\item S2: he can help you and he wants to help you
		\item S3: he wants a computer and a computer for you
	\end{itemize}
\end{exampleblock}

\begin{center}
	\begin{tabular}{lllllllllll}
	\hline\hline
	& a & and & can & computer & for & he & help & to & wants & you \\
	\hline
	S1 & 1 & 0 & 1 & 1 & 0 & 0 & 1 & 0 & 0 & 1\\
	S2 & 0 & 1 & 1 & 0 & 0 & 2 & 2 & 1 & 1 & 2\\
	S3 & 2 & 1 & 0 & 2 & 1 & 1 & 0 & 0 & 1 & 1\\
	\hline\hline
\end{tabular}
\end{center}

\end{frame}


\begin{frame}
\frametitle{Word semantics: Vector representation}
\framesubtitle{TF-IDF: Cosine similarity}

\begin{minipage}{.68\textwidth}
\begin{itemize}
	\item Let $a$ and $b$ be two documents represented by two vectors $\overrightarrow{a}$ and $\overrightarrow{b}$ respectively.
	\item The representation can be vocabulary words (TF or TF-IDF) or other document features.
	\item Let these vectors' length be $n$
\end{itemize}
\end{minipage}
\begin{minipage}{.3\textwidth}
\hgraphpage{exp-cos.pdf}
\end{minipage}

\begin{block}{Cosine similarity between two vectors}
	\[
	Cos(\theta) = \frac{\overrightarrow{a} \overrightarrow{b}}{||\overrightarrow{a}||\, ||\overrightarrow{b}||}
	= \frac{\sum_{i=1}^{n} a_i b_i}{\sqrt{\sum_{i=1}^{n} a_i^2} \sqrt{\sum_{i=1}^{n} b_i^2}}
	\]
\end{block}

\end{frame}

\subsection{Term-term}

\begin{frame}
\frametitle{Word semantics: Vector representation}
\framesubtitle{Term-term}

\begin{itemize}
	\item Words can be represented using other vocabulary words based on co-occurrence.
	\item To represent all words of a vocabulary $ V $, a matrix $|V| \times |V|$ must be used.
	\item Each word is represented by $|V|$ words called \keyword{context}.
	\item Co-occurrence can be calculated according to documents, sentences or windows around the word.
	\item An example of a window: \expword{2 words before and 2 after}.
\end{itemize}

\end{frame}

\begin{frame}
\frametitle{Word semantics: Vector representation}
\framesubtitle{Term-term: Example of a window 2-2}

\begin{exampleblock}{Example of some sentences}
	\begin{itemize}
		\item S1: a computer can help you
		\item S2: he can help you and he wants to help you
		\item S3: he wants a computer and a computer for you
	\end{itemize}
\end{exampleblock}

\begin{center}
	\scriptsize
	\begin{tabular}{lllllllllll}
		\hline\hline
		& a & and & can & computer & for & he & help & to & wants & you \\
		\hline
		a & 0 & 2 & 1 & 4 & 1 & 1 & 0 & 0 & 1 & 0\\
		and & 2 & 0 & 0 & 2 & 0 & 1 & 1 & 0 & 1 & 1\\
		can & 1 & 0 & 0 & 1 & 0 & 1 & 2 & 0 & 0 & 2\\
		computer & 4 & 2 & 1 & 0 & 1 & 0 & 1 & 0 & 1 & 1\\
		for & 1 & 0 & 0 & 1 & 0 & 0 & 0 & 0 & 0 & 1\\
		he & 1 & 1 & 1 & 0 & 0 & 0 & 1 & 1 & 2 & 1\\
		help & 0 & 1 & 2 & 1 & 0 & 1 & 0 & 1 & 1 & 3\\
		to & 0 & 0 & 0 & 0 & 0 & 1 & 1 & 0 & 1 & 1\\
		wants & 1 & 1 & 0 & 1 & 0 & 2 & 1 & 1 & 0 & 0\\
		you & 0 & 1 & 2 & 1 & 1 & 1 & 3 & 1 & 0 & 0\\
		\hline\hline
	\end{tabular}
\end{center}

\end{frame}

\begin{frame}
\frametitle{Word semantics: Vector representation}
\framesubtitle{Term-term: Cosine similarity}

\begin{itemize}
	\item Similarity between two words can be calculated.
	\item Using cosine similarity (previously seen) or other similarity measures such as: Euclidean distance and dot product.
\end{itemize}

\begin{exampleblock}{Example of cosine similarity between ``computer" and ``you"}
	\fontsize{6}{16}\selectfont\bfseries\boldmath
	\begin{align*}
	& \text{Cos(``computer", ``you")} \\ 
	& = \frac{4*0 + 2*1 + 1*2 + 0*1 + 1*1 + 0*1 + 1*3 + 0*1 + 1*0 + 1*0}{\sqrt{4^2 + 2^2 + 1^2 + 0^2 + 1^2 + 0^2 + 1^2 + 0^2 + 1^2 + 1^2} \sqrt{0^2 + 1^2 + 2^2 + 1^2 + 1^2 + 1^2 + 3^2 + 1^2 + 0^2 + 0^2}}\\
	& = \frac{8}{\sqrt{25}\sqrt{18}}\\
	& = \frac{8}{3*5*\sqrt{2}}\\
	\end{align*}
\end{exampleblock}

\end{frame}

\subsection{Latent semantic analysis (LSA)}

\begin{frame}
\frametitle{Word semantics: Vector representation}
\framesubtitle{Latent semantic analysis (LSA)}

\begin{itemize}
	\item Term-document matrix $X[N, M]$ has ...
	\begin{itemize}
		\item a big dimension;
		\item a lot of zeros (see TF example).
	\end{itemize}
	\item $L$ anonymous concepts can be used to represent ...
	\begin{itemize}
		\item terms: $T[N, L]$;
		\item documents: $D[M, L]$.
	\end{itemize}
	\item $L \le \min(N, M)$
	\item $X$ can be decomposed using \keyword{Singular value decomposition (SVD)}
	\begin{itemize}
		\item $X = T \times S \times D^\top$
	\end{itemize}
\end{itemize}

\end{frame}


\begin{frame}
	\frametitle{Word semantics: Vector representation}
	\framesubtitle{LSA: SVD formulation}
	
	\begin{itemize}
		\item $T^\top T = \mathbb{I}_{L \times L}$ 
		\item $D^\top D = \mathbb{I}_{L \times L}$ 
		\item $S$: concepts weights with a decreasing order $s_{11} \ge ... \ge s_{LL} > 0$
		\item SVD can be computed using methods such as \keyword{Lanczos algorithm} and \keyword{QR algorithm}
	\end{itemize}
	
	\begin{block}{SVD of term-document matrix}
		\scriptsize\bfseries
		\[
		\overbrace{
			\begin{bmatrix}
			x_{11} & \ldots & \ldots & \ldots & x_{1M} \\ 
			\vdots & \ddots & \ddots & \ddots &\vdots \\
			\vdots & \ddots & \ddots & \ddots &\vdots \\
			x_{N1} & \ldots & \ldots & \ldots & x_{NM} \\ 
			\end{bmatrix}
		}^{X \text{ : term-document}}
		=
		\overbrace{
			\left[
			\begin{bmatrix}
			t_{11} \\ 
			\vdots \\
			\vdots \\
			t_{N1} \\ 
			\end{bmatrix}
			\begin{matrix}
			\ldots \\ 
			\end{matrix}
			\begin{bmatrix}
			t_{1L} \\ 
			\vdots \\
			\vdots \\
			t_{NL} \\ 
			\end{bmatrix}
			\right]
		}^{T \text{ : term-concept}}
		\times 
		\overbrace{
			\begin{bmatrix}
			s_{11} & \ldots & 0 \\
			0 & \ddots & 0 \\
			0 & \ldots & s_{LL} \\
			\end{bmatrix}
		}^{S \text{ : concept-concept}}
		\times 
		\overbrace{
			\begin{bmatrix}
			\begin{bmatrix}
			d_{11} & \ldots & \ldots & \ldots & d_{1M} \\
			\end{bmatrix}\\
			\vdots \\
			\begin{bmatrix}
			d_{L1} & \ldots & \ldots & \ldots & d_{LM} \\
			\end{bmatrix}\\
			\end{bmatrix}
		}^{D^\top \text{ : concept-document}}
		\]
		
	\end{block}
	
\end{frame}


%===================================================================================
\section{Word embedding}
%===================================================================================

\begin{frame}
\frametitle{Word semantics}
\framesubtitle{Word embedding}

\begin{itemize}
	\item document-term and term-term representations occupy a large memory space.
	\item They are difficult to manage. 
	\item Using \keyword{LSA}, a word's representation is compacted to a much smaller vector of real numbers.
	\item This is called: \keyword{Word embedding}.
	\item In this course, we will focus on neural networks based word embedding.
	\item This technique allows learning some relations between words.
\end{itemize}

\end{frame}


\begin{frame}
\frametitle{Word semantics: Word embedding}
\framesubtitle{Some humor}
\begin{center}
		\vgraphpage{humor/humor-embeddings.jpeg}	
\end{center}
\end{frame}


\subsection{Word2vec}

\begin{frame}
\frametitle{Word semantics: Word embedding}
\framesubtitle{Word2vec}

\begin{itemize}
	\item Using term-term representation, we saw the \keyword{context} notion:
	\begin{itemize}
		\item A word is represented based on the neighboring words.
		\item The vector represents co-occurrence frequencies.
		\item The vector's size is that of the vocabulary (very large).
		\item A word can have co-occurrences with few words (Too much zeros).
	\end{itemize}
	\item Word2vec is a tool provided by Google. It implements two methods of \keyword{Word embedding} \cite{2013-mikolov-al}:
	\begin{itemize}
		\item \optword{CBOW}: Continuous Bag-of-Words;
		\item \optword{Skip-gram}: Continuous Skip-gram.
	\end{itemize}
	\item Encoding words using a small vector (50-1000) based on context (using an encoder-decoder).
\end{itemize}

\end{frame}

\begin{frame}
	\frametitle{Word semantics: Word embedding}
	\framesubtitle{Word2vec: CBOW}
\begin{minipage}{.63\textwidth}
	\begin{itemize}
		\item Infer a word $w_i$ given a context.
		\item In this example, the context is 2 words before and 2 after.
		\item Let $T$ be the number of words in the corpus.
	\end{itemize}
	\begin{block}{Loss function}
		\[%
		J(\theta) = \frac{-1}{T} \sum_{i=1}^{T} \log p(w_i |w_{i-2}, w_{i-1}, w_{i+1}, w_{i+2})
		\]
	\end{block}
\end{minipage}
\begin{minipage}{.36\textwidth}
	\hgraphpage{word2vec-cbow.pdf}
\end{minipage}
	
\end{frame}

\begin{frame}
	\frametitle{Word semantics: Word embedding}
	\framesubtitle{Word2vec: Skip-gram}
	\begin{minipage}{.58\textwidth}
		\begin{itemize}
			\item Infer a context given a word $w_i$.
			\item In this example, the context is 2 words before and 2 after.
			\item Let $T$ be the number of words in the corpus.
		\end{itemize}
		\begin{block}{Loss function}
			\[%
			J(\theta) = \frac{-1}{T} \sum_{i=1}^{T} \sum_{j= i-2; j \ne i}^{i+2} \log p(w_j |w_i)
			\]
		\end{block}
	\end{minipage}
	\begin{minipage}{.4\textwidth}
		\hgraphpage{word2vec-skip.pdf}
	\end{minipage}
	
\end{frame}

\subsection{GloVe}

\begin{frame}
\frametitle{Word semantics: Word embedding}
\framesubtitle{GloVe}

\begin{itemize}
	\item \keyword{GloVe}: Global Vectors.
	\item A method developed at Stanford \cite{2014-pennington-al}.
	\item Trying to exploit both approaches: term-term matrix (such as in LSA) and learning from context (such as in CBOW)
	\item $X[V, V]$ is a term-term matrix where $X_{ij}$ is the number of word $w_j$'s occurrence in the context of the word $w_i$, and $X_i$ is the number of word $w_i$'s occurrences in the corpus.
	\item $P_{ij} = \frac{X_{ij}}{X_i}$ is occurrence probability of $w_j$ in the context of $w_i$.
\end{itemize}

\end{frame}

\begin{frame}
\frametitle{Word semantics: Word embedding}
\framesubtitle{GloVe: Motivation}
	
\begin{minipage}{.5\textwidth}
	\begin{itemize}
		\item We want to find the relation between $w_i = ice$ and $w_j = steam$.
		\item $P_{ik},\ P_{jk}$: probabilities of a word $w_k$ given these two words.
		\item We calculate the ratio: $\frac{P_{ik}}{P_{jk}}$.
	\end{itemize}
\end{minipage}
\begin{minipage}{.48\textwidth}
		\fontsize{6}{10}\selectfont\bfseries
		\begin{tblr}{
				colspec = {p{1.2cm}@{\hskip3pt}p{1cm}@{\hskip3pt}p{1cm}@{\hskip3pt}p{1cm}@{\hskip3pt}p{1.2cm}@{\hskip0pt}},
				row{odd} = {lightblue},
				row{even} = {lightyellow},
				row{1} = {darkblue},
				colsep = 0pt,
				rowsep= 1pt,
			} 
%			\hline\hline
			
			\textcolor{white}{\textbf{Probability and Ratio}} & \textcolor{white}{\textbf{k = solid}} & \textcolor{white}{\textbf{k = gas}} & \textcolor{white}{\textbf{k = water}} & \textcolor{white}{\textbf{k = fashion}} \\
%			\hline
			P(k\textbar ice) & 1.9 * 10\textsuperscript{-4} & 6.6 * 10\textsuperscript{-5} & 3.0 * 10\textsuperscript{-3} & 1.7 * 10\textsuperscript{-5} \\
			P(k\textbar steam) & 2.2 * 10\textsuperscript{-5} & 7.8 * 10\textsuperscript{-4} & 2.2 * 10\textsuperscript{-3} & 1.8 * 10\textsuperscript{-5} \\
%			\hline
			\boldmath\footnotesize$\frac{P(k|ice)}{P(k|steam)}$ & 8.9 & 8.5 * 10\textsuperscript{-2} & 1.36 & 0.96 \\
%			\hline\hline
		\end{tblr}
%		\caption{Exemple de probabilités et un ratio \cite{2014-pennington-al}}
%	\end{figure}
\end{minipage}
	
\begin{itemize}
	\item Ratio's value follows relations $R(w_k, w_i)$ and $R(w_k, w_j)$:
	\begin{itemize}
		\item If $R(w_k, w_i) \wedge \neg R(w_k, w_j)$, the ratio will be large. E.g. \expword{solid}.
		\item If $\neg R(w_k, w_i) \wedge R(w_k, w_j)$, the ratio will be small. E.g. \expword{gas}.
		\item If $R(w_k, w_i) \wedge R(w_k, w_j)$, the ratio will be close to $1$. E.g. \expword{water}.
		\item If $\neg R(w_k, w_i) \wedge \neg R(w_k, w_j)$, the ratio will be close to $1$. E.g. \expword{fashion}.
	\end{itemize}
	\item We want to train a function $F$ on $w_i$, $w_j$ and $\tilde{w_k}$ representations.
	\[F(w_i, w_j, \tilde{w_k}) = \frac{P_{ik}}{P_{jk}}\]
\end{itemize}
	
\end{frame}

\begin{frame}
\frametitle{Word semantics: Word embedding}
\framesubtitle{GloVe: Formulation (intuition)}
	
\begin{itemize}
	\item We want to restrict the function $F$ using substraction:
	\[F(w_i - w_j, \tilde{w_k}) = \frac{P_{ik}}{P_{jk}}\]
	
	\item We want to transform the argument to a scalar using dot product:
	\[F((w_i - w_j)^\top \tilde{w_k}) = \frac{P_{ik}}{P_{jk}}\]
	
	\item We want: $w \leftrightarrow \tilde{w}$ and $X \leftrightarrow X^\top$. 
	So, $F$ must be a homomorphism between $(\mathbb{R}, +)$ and $(\mathbb{R}_{>0}, \times)$:
	\[F((w_i - w_j)^\top \tilde{w_k}) = \frac{F(w_i^\top \tilde{w_k})}{F(w_j^\top \tilde{w_k})}\]
	
	\item So, $F(w_i^\top \tilde{w_k}) = P_{ik} = \frac{X_{ik}}{X_i}$
	
	\item The solution is $F = exp$. So $w_i^\top \tilde{w_k} = \log X_{ik} - \log X_i$
\end{itemize}
	
\end{frame}

\begin{frame}
\frametitle{Word semantics: Word embedding}
\framesubtitle{GloVe: Formulation (Final version)}
	
\begin{itemize}
	\item $log X_i$ is independent from $k$. 
	So, we can train a bias $b_i$ over $w_i$. 
	To respect symmetry, another bias $\tilde{b_k}$ must be trained over $\tilde{w_k}$.
	\item Loss function $J$ is \keyword{least squares method}.
	\begin{itemize}
		\item $J$ must not weight co-occurrences alike: rare co-occurrences must have less impact on $J$.
		\item In this case, $J$ is weighted using a function over $X_{ij}$.
	\end{itemize}
\end{itemize}

\vskip-2pt\begin{block}{GloVe Formulation}\vskip-2pt
	\[w_i^\top \tilde{w_j} + b_i + \tilde{b_j} = \log X_{ij} \]
	\[J(\theta) = \sum_{i=1}^{V} \sum_{j=1}^{V} f(X_{ij}) (w_i^\top \tilde{w_j} + b_i + \tilde{b_j} - \log X_{ij})^2\]
	\vskip-2pt\[f(x) = \begin{cases}
	\frac{x}{x_{max}} & \text{if } x < x_{max} \\
	1 & \text{ otherwise}
	\end{cases}\]
\end{block}
	
\end{frame}

\begin{frame}
\frametitle{Word semantics: Word embedding}
\framesubtitle{GloVe: Architecture}
	
\begin{center}
	\vgraphpage{glove.pdf}
\end{center}
	
\end{frame}

\subsection{Contextual embedding}

\begin{frame}
\frametitle{Word semantics: Word embedding}
\framesubtitle{Contextual embedding}

\begin{exampleblock}{Example of polysemy (word with many senses)}
	\begin{itemize}
		\item My cat is always chasing a \expword{mouse}. (rodent, animal)
		\item The computer's \expword{mouse} is a pointing device. (electronic device)
		\item You are such a \expword{mouse}! (timid person)
	\end{itemize}
\end{exampleblock}

\begin{itemize}
	\item Can Word2vec and GloVe reflect relations between \expword{mouse} and the words cat, computer and timid? How?
	\item Can they distinguish different senses of the word \expword{mouse}? Why?
\end{itemize}

\end{frame}

\begin{frame}
\frametitle{Word semantics: Word embedding}
\framesubtitle{Contextual embedding: ELMo}

\begin{minipage}{.65\textwidth}
\begin{itemize}
	\item \keyword{ELMo}: Embeddings from Language Models.
	\item \url{http://allennlp.org/elmo}
	\item \cite{2018-peters-al}
	\item Representation Characteristics:
	\begin{itemize}
		\item Contextual: takes into account the whole context (sentence).
		\item Deep: combines all intermediate representations of neural network.
		\item Based on characters: allows morphological characteristics and out-of-vocabulary words.
	\end{itemize}
\end{itemize}
\end{minipage}
\begin{minipage}{.33\textwidth}
	\vspace{1.8cm}
	\hgraphpage{Elmo-img.jpg}
\end{minipage}
	
\end{frame}


\begin{frame}
\frametitle{Word semantics: Word embedding}
\framesubtitle{Contextual embedding: ELMo (Architecture)}
	
	\begin{center}
		\vgraphpage{elmo-arch.pdf}
	\end{center}
	
\end{frame}

\begin{frame}
\frametitle{Word semantics: Word embedding}
\framesubtitle{Contextual embedding: ELMo (Description)}

\begin{itemize}
	\item Given a word $w_k$, its representation $x_k$ is calculated based on its characters (see \cite{2015-kim-al}).
	\item $L$ layers of bidirectional LSTM are used. 
	\item For each sentence of $N$ words, maximize\\ 
	\begin{center}
		$\sum_{k=1}^{N} 
	\log P(w_k | w_1,\ldots,w_{k-1}; \Theta_x, \overrightarrow{\Theta}_{LSTM}, \Theta_s)
	+
	\log P(w_k | w_{k+1},\ldots,w_{N}; \Theta_x, \overleftarrow{\Theta}_{LSTM}, \Theta_s)
	$
	\end{center}
	
	\item $R_k = \{x_k^{LM}, \overrightarrow{h}_{LM}^{k, j}, \overleftarrow{h}_{LM}^{k, j} | j= 1 \ldots L \}
	= \{h_{LM}^{k, j} | j= 0 \ldots L \}
	$
	
	\item To integrate ELMo into a $task$, additional parameters $\Theta^{task}$ must be trained to combine intermediate words representations into just one which is task-related:\\
	\begin{center}
		$ELMo_k^{task} = E(R_k; \Theta^{task}) = \gamma^{task} \sum_{j=0}^{L} \theta_j^{task} h_{LM}^{k, j}$
	\end{center}
\end{itemize}

\end{frame}

\begin{frame}
\frametitle{Word semantics: Word embedding}
\framesubtitle{Contextual embedding: ELMo (Some humor)}

\begin{center}
	\vgraphpage{humor/humor-elmo.png}
\end{center}
	
\end{frame}

\begin{frame}
\frametitle{Word semantics: Word embedding}
\framesubtitle{Contextual embedding: BERT}

\begin{minipage}{.71\textwidth}
	\begin{itemize}
		\item \keyword{BERT}: Bidirectional Encoder Representations from Transformers.
		\item \url{https://github.com/google-research/bert}
		\item \cite{2019-devlin-al}
		\item Representation Characteristics:
		\begin{itemize}
			\item Contextual: takes into account the whole context (sentence).
			\item Totally bidirectional: considers direct relations with all sentence's words. 
			\item Based on tokens: words are split into radical + suffixes.
			\item Transfer learning
		\end{itemize}
	\end{itemize}
\end{minipage}
\begin{minipage}{.27\textwidth}
	\vspace{1.1cm}
	\hgraphpage{Bert-img.png}
\end{minipage}
	
\end{frame}

\begin{frame}
\frametitle{Word semantics: Word embedding}
\framesubtitle{Contextual embedding: BERT (Architecture)}
	
	\begin{center}
		\vgraphpage{bert-arch.pdf}
	\end{center}
	
\end{frame}

\begin{frame}
\frametitle{Word semantics: Word embedding}
\framesubtitle{Contextual embedding: BERT (Input)}

\begin{itemize}
	\item Text is tokenized using \keyword{Wordpiece} (see \cite{2016-wu-al}).
	\item The input is a sequence of tokens of maximum length $T = 512$:
	\begin{itemize}
		\item The first token is a special token ``\keyword{[CLS]}" used for classification.
		\item The rest is the words of sentences separated by a token ``\keyword{[SEP]}".
	\end{itemize}
	\item Each token is represented by 3 embeddings of size $N$:
	\begin{itemize}
		\item \optword{Token embedding}: MLP transforming vocabulary OneHot representation of size $V$ into a vector of size $N$.
		\item \optword{Position embedding}: MLP transforming position OneHot representation of size $T$ into a vector of size $N$.
		\item \optword{Segment embedding}: MLP transforming segment OneHot representation of size $2$ (sentence1 or sentence2) into a vector of size $N$.
	\end{itemize}
\end{itemize}
	
\end{frame}

\begin{frame}
\frametitle{Word semantics: Word embedding}
\framesubtitle{Contextual embedding: BERT (Pre-training)}
	
\begin{itemize}
	\item The model is based on \keyword{Transformers}: The encoder part (see \cite{2017-vaswani-al})
	\item It is trained on two different tasks.
	\item \optword{Masked language model}:
	\begin{itemize}
		\item randomly masking 15\% of tokens and trying to infer them.
		\item using a special token \keyword{[MASK]}.
		\item this token will not appear in fine-tuning step.
		This is why just 80\% of replacements use this token.
		\item 10\% with a random token and 10\% without change. 
	\end{itemize}
	\item \optword{Next sentence prediction}:
	\begin{itemize}
		\item Given two sentences in input, predict if the second follow the first. 
		\item For classification purpose, the output of token ``\keyword{[CLS]}" ($CLS \in \{IsNext, NotNext\}$) is used.
	\end{itemize}
	\item Once the model is trained, it can be fine-tuned on another task.
\end{itemize}
	
\end{frame}

\begin{frame}
\frametitle{Word semantics: Word embedding}
\framesubtitle{Contextual embedding: BERT (Fine-tuning)}

	\begin{center}
		\vgraphpage{bert-tasks.pdf}
	\end{center}
	
\end{frame}

\begin{frame}
\frametitle{Word semantics: Word embedding}
\framesubtitle{Contextual embedding: BERT (Some humor)}

	\begin{center}
		\vgraphpage{humor/humor-bert.jpg}
	\end{center}
	
\end{frame}


\subsection{Model evaluation}

\begin{frame}
\frametitle{Word semantics: Word embedding}
\framesubtitle{Model evaluation: Intrinsic evaluation}

\begin{itemize}
	\item  WordSimilarity-353 Test Collection \cite{2002-finkelstein-al}:
	\begin{itemize}
		\item Manually annotated similarities between words (0 et 10).
		\item To test a model's performance, Spearman correlation is used between model's cosine similarities and manual similarities.
	\end{itemize}
	
	\item SimLex-999 \cite{2015-hill-al}:
	\begin{itemize}
		\item The same as the past.
		\item Testing similarity (\expword{coast, shore}) but not association (\expword{clothes, closet}).
	\end{itemize}

	\item Words analogies \cite{2013-mikolov-al2}:
	\begin{itemize}
		\item Dataset is formatted as $(w_{i1}:w_{j1} :: w_{i2}:w_{j2})$.
		\item Test the capacity of embeddings to represent the relations: $w_{j2} = w_{i1} - w_{i2} + w_{j1}$
		\item E.g. \expword{(King:Queen :: Man:Woman) \textrightarrow King - Man + Woman = Queen}
	\end{itemize}

\end{itemize}
	
\end{frame}

\begin{frame}
\frametitle{Word semantics: Word embedding}
\framesubtitle{Model evaluation: Extrinsic evaluation and bias}
	
\begin{itemize}
	\item  \optword{Extrinsic evaluation}: Evaluate the models based on a task:
	\begin{itemize}
		\item Evaluate a new model with respect to a known one according to a task.
		\item E.g. GloVe surpassed LSA in NER task \cite{2014-pennington-al}.
	\end{itemize}
	\item \optword{Bias}: Models can learn biased analogies.
	\begin{itemize}
		\item Learn stereotypes such as \expword{she} with \expword{homemaker, nurse, receptionist} and \expword{he} with \expword{maestro, skipper, protege} \cite{2017-caliskan-al}.
		\item This can affect some tasks' performance. 
		\item E.g. \expword{Anaphora resolution can fail to link the pronoun ``she" with ``doctor"}.
	\end{itemize}
\end{itemize}
	
\end{frame}

%===================================================================================
\section{Word Sense Disambiguation (WSD)}
%===================================================================================

\begin{frame}
\frametitle{Word semantics}
\framesubtitle{Word Sense Disambiguation (WSD)}

\begin{itemize}
	\item Selecting the correct meaning of a word.
	\item Useful for multiple tasks:
	\begin{itemize}
		\item Parsing.
		
		\expword{I \underline{fish} in the river} (Verb or noun?)
		
		\expword{The \underline{fish} was too big} (Verb or noun?)
		
		\item Machine translation.
		
		\expword{I withdrew money from the \underline{bank}} (``sloping land" or ``coin bank"?)
		
		\expword{I fish on the \underline{bank}} (``sloping land" or ``coin bank"?)
	\end{itemize}
\end{itemize}

\end{frame}

\subsection{Knowledge bases based}

\begin{frame}
\frametitle{Word semantics: WSD}
\framesubtitle{Knowledge bases based: Lesk algorithm}

\begin{block}{Lesk algorithm}
	\footnotesize\vspace{-3pt}
	\begin{algorithm}[H]
		\KwData{a word $w$; a sentence $s$ containing $w$}
		\KwResult{The sens of $w$ in $s$}
		
		best\_sens \textleftarrow most frequent sens of $w$\;
		superposition\_max \textleftarrow 0\;
		context \textleftarrow set of words of $s$\; 
		
		\ForEach{sens $w_i$ of $w$}{ 
			signature \textleftarrow set of words in \textbf{gloss} and examples of sens $w_i$\;
			superposition \textleftarrow number of common words between \textbf{context} and \textbf{signature}\;
			\If{superposition $>$ superposition\_max}{
				superposition\_max \textleftarrow superposition\;
				best\_sens \textleftarrow $w_i$\;
			}
		}
		
		\Return best\_sens \;
%		\vspace{-3pt}
	\end{algorithm}
\end{block}


\end{frame}

\begin{frame}
\frametitle{Word semantics: WSD}
\framesubtitle{Knowledge bases based: Graph-based}
	
\begin{itemize}
	\item Example: \url{http://babelfy.org/} \cite{2014-moro-al}
	\item \optword{Step 1: Semantic signatures construction}
	\begin{itemize}
		\item Using a semantic network, attribute a weight for each arc linking two concepts based on the number of triangles linking them.
		\item Calculate a conditional probability of a concept given another based on these weights.
		\item Minimize this graph using ``\keyword{Random walk with restart}".
	\end{itemize}
	\item \optword{Step 2: Candidates identification}
	\begin{itemize}
		\item Apply PoS tagging on the input text.
		\item Extract all possible senses of words and expressions.
	\end{itemize}
	\item \optword{Step 3: Candidates disambiguation}
	\begin{itemize}
		\item Construct a graph based on semantic signature and candidates.
		\item Create a sub-graph by eliminating weak links.
	\end{itemize}
\end{itemize}

\end{frame}

\subsection{ML-based}

\begin{frame}
\frametitle{Word semantics: WSD}
\framesubtitle{ML-based: Supervised ML}
	
\begin{itemize}
	\item Using an annotated corpus (E.g. \expword{SemCor})
	\item Each word is followed by the number of the sens in a lexical base ( such as WordNet).
	\item E.g. \expword{You will find9 that avocado1 is1 unlike1 other1 fruit1 you have ever1 tasted2}
	\item The output of the model is a One-Hot encoded vector over all possible senses. 
	\item In addition to the word, the words surrounding it can be taken as features.
	\item Disambiguation of all words in a text is similar to PoS tagging task.
\end{itemize}
	
\end{frame}

\begin{frame}
\frametitle{Word semantics: WSD}
\framesubtitle{ML-based: Contextual embedding}

\begin{minipage}{.68\textwidth}
\begin{itemize}
	\item 1-nearest-neighbor algorithm
	\item Using contextual embedding (ELMo, BERT, etc.)
	\item Train a model on an annotated corpus to get each token's embedding.
	\item To get a sens embedding $v_s$, calculate the average of all words embeddings $c_i$ which has the same sens.
	\[ v_s = \frac{1}{n} \sum_{i=1}^{n} c_i \] 
\end{itemize}
\end{minipage}
\begin{minipage}{.3\textwidth}
	\begin{figure}
		\hgraphpage{exp-wsd-nn.pdf}
		\caption{Example of WSD with 1-nearest-neighbor \cite{2019-jurafsky-martin}}
	\end{figure}
\end{minipage}

\end{frame}

\insertbibliography{NLP07}{*}

\end{document}

